\documentclass[11pt]{article}

% ---------- packages ----------
\usepackage[a4paper,margin=1in]{geometry}
\usepackage{amsmath,amssymb,amsthm}
\usepackage{mathtools}
\usepackage{bm}
\usepackage{siunitx}
\usepackage{booktabs}
\usepackage{enumitem}
\usepackage{tikz}
\usepackage{pgfplots}
\pgfplotsset{compat=1.18}
\usepackage[colorlinks=true,linkcolor=black,citecolor=black,urlcolor=blue]{hyperref}

% ---------- theorem-like environments ----------
\theoremstyle{remark}
\newtheorem*{remark}{Remark}

% ---------- handy macros ----------
\newcommand{\dd}{\mathrm{d}}
\newcommand{\artanh}{\operatorname{artanh}}
\newcommand{\arsinh}{\operatorname{arsinh}}
\DeclareMathOperator{\sgn}{sgn}

\title{\textbf{Static Shape of a Flexible Cable Hanging\\Between Two Pinned Points}}
\author{}
\date{\today}

\begin{document}
\maketitle

\begin{abstract}
\noindent
A uniform, perfectly flexible, inextensible cable is pinned at two arbitrary points in
a vertical plane and loaded only by its own weight. We derive, from a free-body balance
of an infinitesimal element, the governing ordinary differential equation and integrate
it in closed form. The equilibrium shape is a \emph{catenary}, $y=y_0+a\cosh\!\big((x-x_0)/a\big)$.
We then solve the boundary-value problem posed by the two endpoints and the prescribed
arc length, reducing the three matching conditions to a single transcendental equation
for the catenary parameter~$a$, give the existence criterion, and note the parabolic
small-sag limit. A closing section records the analytical value of the result as a
ground-truth benchmark for numerical cable models.
\end{abstract}

% =====================================================================
\section{Statement and assumptions}
% =====================================================================
A cable hangs in a vertical plane, fixed at two points, under its own weight. We seek its
shape $y(x)$. The model rests on three assumptions.

\begin{enumerate}[label=(A\arabic*),leftmargin=*]
  \item \textbf{Perfectly flexible.} The cable carries no bending moment, so the internal
        force---the tension $T$---is everywhere \emph{tangent} to the curve. This is the
        property that distinguishes a cable from a rod (a rod would add a bending term
        $EI$, raising the order of the governing equation; see Section~\ref{sec:validation}).
  \item \textbf{Inextensible and uniform.} The total arc length is a fixed constant $L$,
        and the weight per unit \emph{arc length} is a constant $w=\lambda g$, with
        $\lambda$ the linear mass density and $g$ gravitational acceleration.
  \item \textbf{Gravity only.} The sole distributed load is the vertical weight; the only
        other forces are the reactions at the two pins.
\end{enumerate}

\noindent
Coordinates: $x$ horizontal, $y$ vertical (positive up). Let $s$ be arc length measured
along the cable and $\theta(s)$ the angle the tangent makes with the horizontal, so that
$\tan\theta = \dd y/\dd x$.

% =====================================================================
\section{Free-body diagram of a cable element}
% =====================================================================
Isolate an infinitesimal element of cable between arc lengths $s$ and $s+\dd s$.
Three forces act on it:
\begin{itemize}[leftmargin=*]
  \item the tension $T(s)$ at the lower cut, directed tangentially ``back-and-down'';
  \item the tension $T(s+\dd s)$ at the upper cut, directed tangentially ``forward-and-up'';
  \item the weight $w\,\dd s$, directed straight down.
\end{itemize}
Resolve the tension into horizontal and vertical components,
\begin{equation}
  H \equiv T\cos\theta, \qquad V \equiv T\sin\theta .
\end{equation}
Static equilibrium of the element yields one scalar equation per direction:
\begin{equation}
  \underbrace{\dd\!\left(T\cos\theta\right)=0}_{\text{horizontal}},
  \qquad
  \underbrace{\dd\!\left(T\sin\theta\right)=w\,\dd s}_{\text{vertical}} .
  \label{eq:element}
\end{equation}

% =====================================================================
\section{The two conservation facts}
% =====================================================================
\paragraph{Horizontal.} The first relation in \eqref{eq:element} integrates immediately:
\begin{equation}
  \boxed{\,H = T\cos\theta = \text{constant}\,}.
  \label{eq:Hconst}
\end{equation}
This is the pivotal fact of the whole problem: because nothing loads the cable
horizontally, the horizontal component of tension is identical at every point.

\paragraph{Vertical.} The second relation gives $\dd V/\dd s = w$. Measure $s$ from the
\emph{vertex}---the lowest point, where the tangent is horizontal, so that $\theta=0$ and
$V=0$ there. Integrating from the vertex,
\begin{equation}
  \boxed{\,V = T\sin\theta = w\,s\,}.
  \label{eq:Vlinear}
\end{equation}
The vertical pull at any point equals the weight of cable hanging between that point and
the vertex.

% =====================================================================
\section{From forces to a slope equation}
% =====================================================================
The slope of the curve is the ratio of the tension components,
\begin{equation}
  \frac{\dd y}{\dd x} = \tan\theta
  = \frac{T\sin\theta}{T\cos\theta}
  = \frac{V}{H}
  = \frac{w\,s}{H}.
\end{equation}
Defining the \emph{catenary parameter}
\begin{equation}
  a \equiv \frac{H}{w} \qquad [\,\text{units of length}\,],
\end{equation}
the slope equation takes the compact form
\begin{equation}
  \frac{\dd y}{\dd x} = \frac{s}{a}.
  \label{eq:slope}
\end{equation}

\begin{remark}
$a$ is the single parameter that fixes the shape: it is the horizontal tension expressed
in ``weight-lengths.'' Large $a$ (taut cable, large $H$) gives a flat curve; small $a$
(slack cable) gives a deep sag.
\end{remark}

% =====================================================================
\section{Closing the equation with the arc-length relation}
% =====================================================================
Equation \eqref{eq:slope} couples $y$ to $s$, and $s$ is itself defined by the curve via
\begin{equation}
  \dd s = \sqrt{1+(y')^2}\,\dd x, \qquad y' \equiv \frac{\dd y}{\dd x}.
\end{equation}
Differentiating \eqref{eq:slope} with respect to $x$ and substituting $\dd s/\dd x$,
\begin{equation}
  \frac{\dd y'}{\dd x} = \frac{1}{a}\frac{\dd s}{\dd x}
  = \frac{1}{a}\sqrt{1+(y')^2}.
  \label{eq:ode}
\end{equation}
This is the governing ODE: a first-order, separable equation for the slope $p\equiv y'$.

% =====================================================================
\section{First integration: the slope}
% =====================================================================
Separating variables in \eqref{eq:ode},
\begin{equation}
  \frac{\dd p}{\sqrt{1+p^2}} = \frac{\dd x}{a}.
\end{equation}
The left-hand side integrates to $\arsinh p$, since
$\tfrac{\dd}{\dd p}\arsinh p = 1/\sqrt{1+p^2}$, giving
\begin{equation}
  \arsinh(p) = \frac{x}{a} + C_1.
\end{equation}
Placing $x=0$ at the vertex, where $p=y'=0$, fixes $C_1=\arsinh(0)=0$. Inverting,
\begin{equation}
  \boxed{\,\frac{\dd y}{\dd x} = \sinh\frac{x}{a}\,}.
  \label{eq:slopesol}
\end{equation}

% =====================================================================
\section{Second integration: the shape}
% =====================================================================
Integrating \eqref{eq:slopesol},
\begin{equation}
  y = \int \sinh\frac{x}{a}\,\dd x = a\cosh\frac{x}{a} + C_2.
\end{equation}
The constant $C_2$ merely places the curve vertically. Writing it together with a
horizontal shift gives the \emph{catenary}:
\begin{equation}
  \boxed{\,y(x) = y_0 + a\cosh\frac{x-x_0}{a}\,},
  \label{eq:catenary}
\end{equation}
whose vertex sits at $\big(x_0,\,y_0+a\big)$. The shape is a hyperbolic cosine---%
\emph{not} a parabola (Section~\ref{sec:parabola}).

\begin{figure}[h]
\centering
\begin{tikzpicture}
\begin{axis}[
    width=12cm, height=6.2cm,
    axis lines=middle,
    xlabel={$x$}, ylabel={$y$},
    xmin=-3.4, xmax=3.4, ymin=0, ymax=6.2,
    xtick=\empty, ytick=\empty,
    clip=false,
    samples=100,
]
% catenary y = a cosh(x/a) with a = 1, vertex at (0,1)
\addplot[thick, domain=-3:2.2]{cosh(x)};
% chord between the two pins
\addplot[dashed, gray] coordinates {(-3,{cosh(-3)}) (2.2,{cosh(2.2)})};
% pins
\addplot[only marks, mark=*, mark size=2pt] coordinates {(-3,{cosh(-3)}) (2.2,{cosh(2.2)})};
\node[anchor=south east] at (axis cs:-3,{cosh(-3)}) {$P_1=(x_1,y_1)$};
\node[anchor=south west] at (axis cs:2.2,{cosh(2.2)}) {$P_2=(x_2,y_2)$};
% vertex
\addplot[only marks, mark=*, mark size=1.5pt, color=black] coordinates {(0,1)};
\node[anchor=north] at (axis cs:0,1) {vertex};
\node[anchor=west] at (axis cs:0.05,3.1) {$y=y_0+a\cosh\!\frac{x-x_0}{a}$};
\end{axis}
\end{tikzpicture}
\caption{A catenary pinned at two points of unequal height. The dashed line is the chord
of length $\sqrt{h^2+v^2}$; the cable's arc length between the pins is $L>\sqrt{h^2+v^2}$.
At the vertex the tension is purely horizontal and equal to $H=wa$.}
\label{fig:catenary}
\end{figure}

% =====================================================================
\section{Arc length and a built-in consistency check}
% =====================================================================
Using \eqref{eq:slopesol} and the identity $1+\sinh^2=\cosh^2$,
\begin{equation}
  s(x) = \int_0^x \sqrt{1+\sinh^2\tfrac{x'}{a}}\;\dd x'
       = \int_0^x \cosh\tfrac{x'}{a}\,\dd x'
       = a\sinh\frac{x}{a}.
  \label{eq:arclen}
\end{equation}
Substituting back into \eqref{eq:slope} returns $y'=s/a=\sinh(x/a)$, in agreement with
\eqref{eq:slopesol}: the derivation closes on itself.

\begin{remark}[Tension distribution]
From \eqref{eq:Hconst} and $\cos\theta = 1/\sqrt{1+\sinh^2} = 1/\cosh\!\frac{x-x_0}{a}$,
the tension is
\[
  T = H\cosh\frac{x-x_0}{a} = w\,(y-y_0).
\]
Tension is proportional to height above the directrix $y=y_0$. The minimum, $H=wa$, occurs
at the vertex; the maximum occurs at the higher support. This is the relation to consult
when checking whether a cable, or a discretized joint chain, will exceed its breaking load.
\end{remark}

% =====================================================================
\section{The boundary-value problem}
% =====================================================================
The catenary \eqref{eq:catenary} carries three unknowns: $a$ (shape), $x_0$ (vertex
abscissa), and $y_0$ (vertical offset). The data supply exactly three conditions---the two
pins and the length:
\begin{equation}
  \text{(i) } y(x_1)=y_1,\qquad
  \text{(ii) } y(x_2)=y_2,\qquad
  \text{(iii) } s(x_2)-s(x_1)=L.
\end{equation}
Introduce the dimensionless endpoint arguments
\begin{equation}
  A\equiv\frac{x_1-x_0}{a},\qquad B\equiv\frac{x_2-x_0}{a},
\end{equation}
so that the three conditions read
\begin{equation}
  y_1 = y_0+a\cosh A,\qquad
  y_2 = y_0+a\cosh B,\qquad
  L = a\,(\sinh B-\sinh A).
\end{equation}

% =====================================================================
\section{Reduction to a single transcendental equation}
% =====================================================================
Let $h\equiv x_2-x_1$ be the horizontal span and $v\equiv y_2-y_1$ the vertical rise.
Subtracting (i) from (ii) eliminates $y_0$:
\begin{equation}
  v = a\,(\cosh B-\cosh A),\qquad\qquad
  L = a\,(\sinh B-\sinh A).
\end{equation}
Apply the sum-to-product identities
\begin{align}
  \cosh B-\cosh A &= 2\sinh\tfrac{B+A}{2}\sinh\tfrac{B-A}{2},\\
  \sinh B-\sinh A &= 2\cosh\tfrac{B+A}{2}\sinh\tfrac{B-A}{2}.
\end{align}
Since $B-A=(x_2-x_1)/a=h/a$, we have $\tfrac{B-A}{2}=\tfrac{h}{2a}$. Writing
$m\equiv\tfrac{A+B}{2}$,
\begin{equation}
  v = 2a\,\sinh m\,\sinh\frac{h}{2a},
  \qquad\qquad
  L = 2a\,\cosh m\,\sinh\frac{h}{2a}.
  \label{eq:vL}
\end{equation}
Combine the two equations in \eqref{eq:vL} in two ways. \emph{Dividing} cancels the
common factor $2a\sinh\frac{h}{2a}$:
\begin{equation}
  \frac{v}{L}=\tanh m
  \quad\Longrightarrow\quad
  m=\artanh\frac{v}{L}.
  \label{eq:m}
\end{equation}
\emph{Squaring and subtracting}, with $\cosh^2 m-\sinh^2 m=1$,
\begin{equation}
  L^2-v^2 = \left(2a\sinh\tfrac{h}{2a}\right)^2,
\end{equation}
which yields the master equation---a single transcendental equation in the lone
unknown $a$:
\begin{equation}
  \boxed{\,\sqrt{L^2-v^2} = 2a\,\sinh\frac{h}{2a}\,}.
  \label{eq:master}
\end{equation}

% =====================================================================
\section{Solution recipe}
% =====================================================================
\begin{enumerate}[label=\textbf{\arabic*.},leftmargin=*]
  \item \textbf{Solve \eqref{eq:master} for $a$.} The right-hand side
        $f(a)=2a\sinh\frac{h}{2a}$ is monotone (Section~\ref{sec:existence}), so Newton's
        method or bisection converges quickly. A serviceable initial guess follows from the
        small-sag relation $f(a)\approx h\big(1+\tfrac{h^2}{24a^2}\big)$, i.e.
        $a_0\approx h\big/\!\sqrt{24\,(\sqrt{L^2-v^2}/h-1)}$.
  \item \textbf{Vertex abscissa $x_0$} from \eqref{eq:m}. Since
        $m=\tfrac{A+B}{2}=\dfrac{x_1+x_2-2x_0}{2a}$,
        \begin{equation}
          x_0 = \frac{x_1+x_2}{2} - a\,\artanh\frac{v}{L}.
        \end{equation}
  \item \textbf{Vertical offset $y_0$} from condition (i):
        \begin{equation}
          y_0 = y_1 - a\cosh\frac{x_1-x_0}{a}.
        \end{equation}
  \item \textbf{Assemble} the shape:
        \begin{equation}
          y(x) = y_0 + a\cosh\frac{x-x_0}{a}.
        \end{equation}
\end{enumerate}

% =====================================================================
\section{Existence and the taut limit}
\label{sec:existence}
% =====================================================================
Consider $f(a)=2a\sinh\frac{h}{2a}$. As $a\to0^+$ the hyperbolic sine dominates and
$f\to\infty$; as $a\to\infty$, $\sinh\varepsilon\approx\varepsilon$ gives
$f\to 2a\cdot\frac{h}{2a}=h$. The function decreases monotonically from $\infty$ to $h$.
Hence \eqref{eq:master} has a unique positive root if and only if
\begin{equation}
  \sqrt{L^2-v^2}>h
  \;\Longleftrightarrow\;
  L>\sqrt{h^2+v^2}.
\end{equation}
The cable must be longer than the chord. When $L$ equals the chord, $a\to\infty$ and the
catenary degenerates to a straight, taut line; $L$ less than the chord is impossible for an
inextensible cable---precisely the expected physical constraint.

% =====================================================================
\section{Why it is a cosh and not a parabola}
\label{sec:parabola}
% =====================================================================
A suspension-bridge main cable \emph{is} parabolic, because there the load is uniform per
unit \emph{horizontal} distance (the deck), giving $V=w_x x$ and $y''=\text{const}$. Here
the load is uniform per unit \emph{arc length} (the cable's own weight), which produces the
$\cosh$. The parabola is the small-sag limit of the catenary: for $|x-x_0|\ll a$,
\begin{equation}
  a\cosh\frac{x-x_0}{a} \approx a + \frac{(x-x_0)^2}{2a},
\end{equation}
the leading term of the Taylor expansion. A shallow catenary therefore looks parabolic,
but the exact self-weight solution is the hyperbolic cosine.

% =====================================================================
\section{The result as a numerical benchmark}
\label{sec:validation}
% =====================================================================
Equation \eqref{eq:catenary} with the constants of the solution recipe is the analytical
ground truth for a static hanging cable, and as such is the canonical check for numerical
rod and cable models:
\begin{itemize}[leftmargin=*]
  \item A \textbf{position-based (XPBD) inextensible rod} should relax to this catenary.
        Residual error indicates insufficient constraint stiffness or too few solver
        iterations, not a modelling defect.
  \item A \textbf{capsule chain with free angular joints} matches the catenary; any
        non-zero joint stiffness appears as a systematic departure from $\cosh$ that this
        solution quantifies exactly.
  \item A \textbf{volumetric FEM cable} will \emph{not} match it, because it absorbs axial
        load as stretch. Its rest shape is the \emph{elastic catenary}, in which the
        arc-length--tension relation acquires a stretching term proportional to $1/(EA)$;
        under identical endpoints it sags more than the inextensible solution. That is the
        appropriate ground truth for the two-way FEM case and is a natural next derivation.
\end{itemize}
A final caveat on scope: the present result assumes zero bending stiffness. A rod with
non-negligible $EI$ is no longer a catenary but a gravito-elastica, governed by a
fourth-order ODE with a boundary layer near each clamped end. For a cable, $EI\approx0$ and
the catenary is exact.

\end{document}
